
\section{BioseqDB SIMD performance comparison}

I perfomed intermediate testing to verify if predictions coming from synthetic results can be correct.

The test used $|T|=500Mbps$ and $|Q|=5Mbps$ sequences  with just one $N=1$ occurence (this means that this is very optimistic case of random scenario)

\begin{figure}[H]
\centering
\begin{tikzpicture}[scale=1.0,transform shape,x=1cm,y=1cm,
  font=\footnotesize,>=Latex]
  \node[anchor=west,font=\bfseries] at (.10,9.55)
    {(a) Chaining};
  \draw[->] (2.50,5.10)--(12.15,5.10) node[right] {ms};
  \foreach \x/\lab in {0/0,20/20,40/40,60/60,80/80,100/100,120/120}
    {
      \draw (2.50+.075*\x,5.00)--(2.50+.075*\x,5.20);
      \node[anchor=north] at (2.50+.075*\x,4.97) {\lab};
    }
  \node[anchor=east] at (2.35,8.55) {B1 scalar};
  \node[anchor=east] at (2.35,7.35) {B1 exact SIMD};
  \node[anchor=east] at (2.35,6.15) {B5 bounded SIMD};
  \draw[fill=chainGray!32,draw=chainGray!90!black]
    (2.50,8.17) rectangle (2.50+.075*115.286,8.93);
  \draw[fill=chainBlue!22,draw=chainBlue!90!black]
    (2.50,6.97) rectangle (2.50+.075*101.885,7.73);
  \draw[fill=chainOrange!25,draw=chainOrange!95!black]
    (2.50,5.77) rectangle (2.50+.075*67.012,6.53);
  \node[anchor=west] at (2.66+.075*115.286,8.55) {115.286};
  \node[anchor=west] at (2.66+.075*101.885,7.35) {101.885};
  \node[anchor=west] at (2.66+.075*67.012,6.15) {67.012};

  \node[anchor=west,font=\bfseries] at (.10,4.15)
    {(b) Backtrack};
  \draw[->] (2.50,-.30)--(10.10,-.30) node[right] {ms};
  \foreach \x/\lab in {0/0,2/2,4/4,6/6,8/8,10/10,12/12}
    {
      \draw (2.50+.58*\x,-.40)--(2.50+.58*\x,-.20);
      \node[anchor=north] at (2.50+.58*\x,-.43) {\lab};
    }
  \node[anchor=east] at (2.35,3.15) {B1 scalar};
  \node[anchor=east] at (2.35,1.95) {B1 exact SIMD};
  \node[anchor=east] at (2.35,.75) {B5 bounded SIMD};
  \draw[fill=chainGray!32,draw=chainGray!90!black]
    (2.50,2.77) rectangle (2.50+.58*11.769,3.53);
  \draw[fill=chainBlue!22,draw=chainBlue!90!black]
    (2.50,1.57) rectangle (2.50+.58*8.549,2.33);
  \draw[fill=chainOrange!25,draw=chainOrange!95!black]
    (2.50,.37) rectangle (2.50+.58*8.345,1.13);
  \node[anchor=west] at (2.65+.58*11.769,3.15) {11.769};
  \node[anchor=west] at (2.65+.58*8.549,1.95) {8.549};
  \node[anchor=west] at (2.65+.58*8.345,.75) {8.345};
\end{tikzpicture}
\caption{Recorded times for the three CPU modes}
\label{fig:chaining-mode-times}
\end{figure}

The figure above depicts only two runs and presents only quick prototyping test measurements to judge if we are going in right direction.
Please note that on bottom chart B1 and B5 are similar, because they use the same backtracking implementation (this is the same algorithm).